\emph{Reduction to quadratic forms.} Under (P) and (C) an admissible density is a polynomial in the field strength $F$ whose coefficients are polynomials in the algebraic data $I$, $X$ (or $V$), formed with wedge products and no background form (a fixed coordinate four-form is not a field of the series). A product of $k$ factors $F$ is a $2k$-form, so a four-form on four-dimensional spacetime is homogeneous of degree two in $F$; two-forms commute under the wedge product, so $L = B(F \wedge F)$ for a symmetric bilinear form $B$ on the values of $F$, depending on the data. A gauge transformation acts on the values by $F \mapsto UF\tilde{U}$ with $U$ a rotor of the gauge group, and on the data by $V \mapsto UV\tilde{U}$; (G) is invariance under these. With $V$ ranging over its orbit $\{Q(V) = -3/\Lambda,\ \ip{V}{I} = 0\}$ under the stabiliser of $I$, the form $B_{V}$ is fixed by $B_{V_{0}}$ at one point $V_{0}$, which must be invariant under the stabiliser of $(I, V_{0})$; conversely the three forms $K$, $H$, $M$ below are polynomial in $V$, so each invariant $B_{V_{0}}$ extends polynomially. The problem is therefore the finite one of computing the invariant symmetric bilinear forms of a Lie algebra of bivectors on a space of bivectors.

\emph{The counts (i) to (iii).} The infinitesimal invariance condition $B([\xi, u], v) + B(u, [\xi, v]) = 0$ for $\xi$ in a spanning set of the algebra is a homogeneous linear system on the coefficients of $B$; the check solves it exactly over the rationals, with the stabilisers computed as the bivectors $\xi$ with $\xi \times I = 0$ (and $\xi \times V = 0$), dimensions $10$ and $6$, at $\Lambda \in \{3, -3, \tfrac97\}$. Results: $\so(2,4)$ on all bivectors, $1$; its parabolic algebra of the line of $X$ (the stabiliser of $X$ with the dilation $\nO \wedge \nInf$, which maps $X$ to itself), $1$; the stabiliser of $X$ alone on all bivectors, $2$, with no form odd under the internal reflection in $e_{1}$ (which fixes $X$); the stabiliser of $X$ on $\Lambda^{2}(X^{\perp})$, which is the stabiliser of $X$ as a space, $2$, exactly one of them internally odd (on this subspace $\epsilon(I, X, F, G)$ depends on $I$ only through $\ip{X}{I}$; it is not the restriction of an invariant form on all bivectors, since those restrict to even forms); the stabiliser of $I$ on $\Lambda^{2}(I^{\perp})$, $1$; the stabiliser of $(I, V)$ on $\Lambda^{2}(I^{\perp})$, $3$; on all bivectors, $6$. These agree with the representation count: $\so(2,4)$, $\so(1,4)$ and $\so(2,3)$ are simple, so their invariant symmetric forms are the multiples of the Killing form, while under the Lorentz algebra $\Lambda^{2}(I^{\perp}) = \Lambda^{2}W \oplus (W \wedge V)$, with two invariant forms on $\Lambda^{2}W$ (the scalar and pseudoscalar contractions) and one on $W \wedge V$, and no cross term since $\Lambda^{2}W$ and $W$ share no irreducible summand. The check confirms that $K$, $H$ and $M$ lie in the three-dimensional space and are independent, that $M$ computed with $X$ in place of $V$ is the same form (because $V - X = -\tfrac{3}{\Lambda}I$ and $I \wedge I = 0$), that the internal reflection $u \mapsto e_{1}ue_{1}^{-1}$, which fixes $I$ and $V$, fixes $K$ and $H$ and reverses $M$, that the odd subspace is exactly the multiples of $M$ and the even one is spanned by $K$ and $H$, and that of the three only $K$ is invariant under the whole stabiliser of $I$.

\emph{Chern-Weil forms are topological.} If $B$ is invariant under the algebra $\mathfrak{g}$ in which $A$ takes values, then $dB(F \wedge F) = 2B(DF \wedge F) = 0$ by the Bianchi identity $DF = 0$, the density is locally $d$ of the Chern-Simons form of $A$, and its Euler-Lagrange derivative vanishes identically (Chern-Weil theory and Chern and Simons 1974, Fact \claimref{synthesis.imported-chern-weil-and-chern-simons-forms}(a), (b)). This gives the topological statements of (i), (ii) and the $K$ part of (iii); for a connection valued in the stabiliser of $X$ it makes both invariant forms on $\Lambda^{2}(X^{\perp})$ Chern-Weil. On the normal tractor connection, which is not valued in the stabiliser of $X$ although its field strength annihilates $X$ when the torsion vanishes, the odd form is not topological: it is the Weyl square (Theorem \claimref{synthesis.frozen-point-integrand-on-tractor-transport-is-weyl-squared}(iii)).

\emph{Parity.} The internal reflection acts on the values only, with the spacetime orientation fixed. The integral of a four-form changes sign under reversal of the spacetime orientation, so under the combined parity (orientation reversal together with the internal reflection, the frame identifying the two) a density odd under the internal reflection is invariant and an internally even one changes sign. Hence $M$ (and its Euler, Einstein-Hilbert and volume parts) is even under the combined parity and $K$, $H$ (Pontryagin and Holst type) are odd, the usual assignment.

\emph{Field equations (iv).} Let $B$ be invariant only under the Lorentz algebra $\mathfrak{h}$ of $(I, V)$ and write $A = \omega + \theta$ with $\omega$ in $\mathfrak{h}$ and $\theta(a) = \hb^{-1}(a) \wedge V/Q(V)$ in the complement $W \wedge V$: in the gauge $V$ constant this is the general connection preserving $I$, the field variable for every field (on Einstein fields it is the connection of Theorem \claimref{synthesis.gtg-tractor-gauge-field-reproduces-tractor-connection}(v)). Then $dB(P \wedge R) = B(D_{\omega}P \wedge R) \pm B(P \wedge D_{\omega}R)$, $D_{A} = D_{\omega} + [\theta, \cdot\,]$, and $D_{\omega}F = -[\theta, F]$ by the Bianchi identity, so for a variation $\delta A$,
\[
\delta B(F \wedge F) = 2\,dB(\delta A \wedge F) + 2B([\theta, \delta A] \wedge F) - 2B(\delta A \wedge [\theta, F]) .
\]
For $H$ and $M$, which see only the $\Lambda^{2}W$ block, and a frame variation $\delta A = \delta\hb^{-1} \wedge V/Q(V)$ in $W \wedge V$, the last term vanishes and $[\theta, \delta A]$ is a nonzero multiple of $\hb^{-1} \wedge \delta\hb^{-1}$ in $\Lambda^{2}W$, so the frame equation is the three-form $B(e \wedge \delta e \wedge F_{W})$ with $e = \hb^{-1}$ and $F_{W} = -R + \tfrac{\Lambda}{3}e \wedge e$ the Lorentz block (the store's sign, vanishing on the de Sitter field). At a point, with the frame the identity and $R$ any of the $36$ curvature two-forms with values in $\Lambda^{2}W$, the check evaluates the sixteen components $E_{B}(d, u) = \sum\epsilon^{d\nu\rho\sigma}B(u \wedge e_{\nu}, F_{W}(\rho, \sigma))$. The first Bianchi identity (the cyclic sum of $R(\rho\sigma) \cdot e_{\tau}$, which is $DT$ in Cartan's structure equations, a labelled reading) cuts out a $20$-dimensional space, the algebraic curvature tensors. For $M$ the map has rank $10$ there (the Weyl kernel of Theorem \claimref{synthesis.macdowell-mansouri-principle-fixes-the-trace-free-part}(ii)) and rank $16$ on all $36$, it vanishes on the de Sitter curvature $R(\rho\sigma) = \tfrac{\Lambda}{3}e_{\rho} \wedge e_{\sigma}$ and not at three times it, and its value at $R = 0$ (the volume term) is nonzero. For $H$ the map is zero on the $20$ algebraic curvatures, rank $16$ on all $36$, and its value at $R = 0$ is zero. So $M$ and $H$ are not topological (their frame equations do not vanish identically), $H$ is fed only by the non-cyclic part, which is the torsion's derivative, and by the connection equation of Fact \claimref{synthesis.imported-einstein-cartan-matter-coupling} the torsion vanishes without spin density. With spinning matter the same Fact makes the torsion proportional to the spin density, so $H$ then changes the field equations (through the spin-spin contact term, the readings recorded in Fact \claimref{synthesis.imported-lovelock-cartan-four-forms-and-holst-term}): without the parity hypothesis, uniqueness of $\kappa M$ is equivalent to a torsion-free coupling. $H$ is the Holst term up to a Pontryagin density: on $\Lambda^{2}W$ it is the scalar contraction, and $H(F_{W}, F_{W}) = H(R, R) - \tfrac{2\Lambda}{3}H(R, e \wedge e)$ because $H(e \wedge e, e \wedge e) = 0$ (Holst 1996, Fact \claimref{synthesis.imported-lovelock-cartan-four-forms-and-holst-term}(b), which also gives that without spin the connection equation with $H$ added still forces zero torsion).

\emph{(v).} By (iii) an admissible action is $\kappa M + \kappa' H + \kappa'' K$, and invariance under the combined parity leaves $\kappa M$; $M = \epsilon(I, V, F, F) = \epsilon(I, X, F, F)$, and $Q(I)^{-1}\kappa M$ is the frozen-point functional of Theorem \claimref{dark-energy.frozen-point-makes-stelle-west-vector}, whose split into Gauss-Bonnet, Einstein-Cartan and cosmological terms is Fact \claimref{dark-energy.imported-macdowell-mansouri-theorem-stelle-west}. The split stated is bilinearity of $M$ in $F_{W} = -R + \tfrac{\Lambda}{3}e \wedge e$ with $Q(I)^{-1} = 3/\Lambda$. No invariance fixes $\kappa$. Any admissible internally odd density is a multiple of $M$, and the Lorentz Euler density $M(R, R)$ is not admissible because $R = -F_{W} + \tfrac{\Lambda}{3}e \wedge e$ involves $e$, a multiple of $DV$, so adding one rescales all three terms together: the Euler term cannot be removed inside the class, and the volume-to-Einstein-Hilbert ratio, hence the field equations' $\Lambda$, is the frozen point's.

\emph{(vi).} If $V$ may also enter through $DV$, then in the gauge $V$ constant $DV = [\theta, V]$ is a nonzero multiple of the soldering form $e$, a $W$-valued one-form ($\ip{V}{DV} = 0$ since $Q(V)$ is constant, $\ip{I}{DV} = 0$ since $I$ is parallel and $\ip{I}{V} = 0$). The additional four-forms are $B'(e \wedge e \wedge F)$ with $B'$ a Lorentz-invariant bilinear form on $\Lambda^{2}W \times \Lambda^{2}(I^{\perp})$, of which the check finds $2$ (the pseudoscalar contraction, Einstein-Hilbert, and the scalar one, Holst), and $e^{4}$, a multiple of the one invariant of $\Lambda^{4}W$ (volume): $3 + 2 + 1 = 6$, comparable with the Lovelock-Cartan classification in four dimensions (Mardones and Zanelli 1991, Fact \claimref{synthesis.imported-lovelock-cartan-four-forms-and-holst-term}(a); their list also has torsion-squared and exact terms, so the comparison is not term by term). Their coefficients are independent, so the Euler term can then be subtracted, $\Lambda$ in the field equations is a free coefficient, and the frozen point fixes only the constraint data. Without the reduction ($F$ any bivector two-form, invariance under the Lorentz algebra of $(I, V)$) the check finds $6$ invariant forms. Finally, under (T), that the special-conformal slot of the tractor gauge field is a function of $(\hb, \Omega)$ and involves no further field, it must be a $W$-valued one-form built from them, covariant, polynomial and inverse-frame-free (this step is proved as clause (i) of Theorem \claimref{synthesis.inverse-frame-free-tractor-gauge-field-preserves-a-frozen-point}, and without (T) an independent slot field also satisfies (P), part (iv) there). Under (P) the covariant forms available are $e$ (degree one), the curvature and torsion (degree two) and their covariant derivatives; with no inverse frame no contraction lowers the degree and no non-constant scalar zero-form can be built, so the slot is linear in $e$ with a constant coefficient, $e \mapsto S(e)$ a Lorentz-equivariant map of $W$. The check finds the Lorentz commutant on $W$ one-dimensional, so $S = \rho\,\mathrm{id}$; then $\Theta(A) \times (\nInf + \rho\nO) = \rho A - S(A) = 0$ (Theorem \claimref{synthesis.gtg-tractor-gauge-field-reproduces-tractor-connection}(ii)), checked exactly, and not for the slot $2\rho\,\mathrm{id}$. The parallel datum is a frozen point with $\Lambda = 6\rho$. The Schouten slot of the normal tractor connection is a contraction of the curvature with the inverse frame, so the normal connection is not admissible.

\emph{Scope.} Proved: the classification within (P), (C), (G) at the stated algebras, and the field-equation ranks at a point. Imported: Chern-Weil closedness (Fact \claimref{synthesis.imported-chern-weil-and-chern-simons-forms}), the absence of torsion without spin (the imported coupling, with the Holst term inert by Fact \claimref{synthesis.imported-lovelock-cartan-four-forms-and-holst-term}(b)), and the value of $\epsilon(I, X, F, F)$ on the normal tractor connection (Theorem \claimref{synthesis.frozen-point-integrand-on-tractor-transport-is-weyl-squared}); the Lovelock-Cartan comparison and the Barbero-Immirzi readings are cited, not used. Not addressed: why there should be a local action at all, and the matter coupling.
